%%%%%%%%%%%%%%%%%%%%%%%%%%%%%%%%%%%%%%%%%%%%%%%%%%%%%%%%%%%%%%%%%%%%%%%%%%%%%%%
% APPENDIX AE: Mechanism Design
% Designing Institutions to Achieve Desired Outcomes
%%%%%%%%%%%%%%%%%%%%%%%%%%%%%%%%%%%%%%%%%%%%%%%%%%%%%%%%%%%%%%%%%%%%%%%%%%%%%%%
\section{Mechanism Design: Engineering Economic Institutions}
\label{app:mechanism-design}

%==============================================================================
% PART 0: INTEGRATION OVERVIEW
%==============================================================================
\subsection{Framework Integration Map}

This appendix demonstrates how mechanism design provides essential foundations 
for the $(C, \Psi, C^*, K, Q)$ framework. Mechanism design is ``reverse game 
theory'': instead of analyzing $C^*$ for a given game, it designs the game 
($\Psi_{mechanism}$) to achieve a desired outcome. This is the science of 
$\Psi$-design---creating institutional rules that align individual incentives 
with social objectives. The 2007 Nobel Prize to Hurwicz, Maskin, and Myerson 
recognized this as foundational for institutional economics.

\begin{center}
\renewcommand{\arraystretch}{1.4}
\begin{tabular}{|l|l|l|}
\hline
\textbf{Framework Element} & \textbf{Mechanism Design Contribution} & \textbf{Key Papers} \\
\hline
$\Psi_{mechanism}$ design & Reverse game theory & Hurwicz 1960, 1972 \\
$C^* = C^*_{truthful}$ & Revelation Principle & Myerson 1979 \\
Implementability & Which $C^*$ achievable? & Maskin 1999 \\
Impossibility & Limits of $\Psi$ design & Gibbard 1973, Myerson-Satterthwaite 1983 \\
Optimal $\Psi$ & Revenue/welfare maximization & Myerson 1981, Vickrey 1961 \\
Contract as $\Psi$ & Moral hazard, adverse selection & Holmström 1979 \\
\hline
\end{tabular}
\end{center}

\paragraph{Mechanism Design's Unique Role.}
Mechanism design contributes uniquely to the framework:
\begin{enumerate}
    \item \textbf{$\Psi$ as design variable}: Institutions are chosen, not given
    \item \textbf{Incentive constraints}: What $C^*$ can be implemented?
    \item \textbf{Revelation Principle}: Truth-telling as benchmark
    \item \textbf{Impossibility theorems}: Fundamental limits on design
    \item \textbf{Optimal mechanisms}: Best $\Psi$ given constraints
\end{enumerate}

%==============================================================================
% PART I: FOUNDATIONS – HURWICZ AND THE DESIGN PROBLEM
%==============================================================================
\subsection{Part I: Foundations – The Design Problem}

\subsubsection{Hurwicz: Mechanism Design as Science}

Leonid Hurwicz (1960, 1972) founded mechanism design:

\begin{equation}
\boxed{
\text{Given desired } f: \Theta \to X, \text{ find mechanism } M \text{ that implements } f
}
\label{eq:hurwicz}
\end{equation}

where $\Theta$ is type space (private information), $X$ is outcome space, and 
$f$ is social choice function.

\paragraph{The Problem.}
\begin{itemize}
    \item Designer wants outcome $f(\theta)$ but doesn't know $\theta$
    \item Agents know their own $\theta_i$ (private information)
    \item Must design rules so agents reveal or act on $\theta$ truthfully
\end{itemize}

\subsubsection{Framework Translation}

Mechanism design is $\Psi$-engineering:
\begin{equation}
\boxed{
\Psi_{mechanism}^* = \arg\max_{\Psi} Q(C^*(\Psi))
}
\label{eq:psi-design}
\end{equation}

Choose institutional rules $\Psi$ to maximize welfare $Q$ at equilibrium $C^*$.

\subsubsection{Components of a Mechanism}

A mechanism $M = (S, g)$ consists of:
\begin{itemize}
    \item \textbf{Strategy space} $S = S_1 \times \cdots \times S_n$: What agents can do
    \item \textbf{Outcome function} $g: S \to X$: How strategies map to outcomes
\end{itemize}

\paragraph{Framework Translation.}
\begin{equation}
\Psi_{mechanism} = (S, g) \quad \text{determines possible } C \text{ and resulting } X
\end{equation}

\subsubsection{Implementation}

A mechanism $M$ \textbf{implements} social choice function $f$ if:
\begin{equation}
g(C^*(M)) = f(\theta) \quad \forall \theta
\end{equation}

At equilibrium, the mechanism produces the desired outcome.

\paragraph{Types of Implementation.}
\begin{center}
\renewcommand{\arraystretch}{1.3}
\begin{tabular}{|l|l|l|}
\hline
\textbf{Type} & \textbf{Equilibrium Concept} & \textbf{Strength} \\
\hline
Dominant strategy & Truth optimal regardless of others & Strongest \\
Nash & Truth optimal given others truthful & Medium \\
Bayesian Nash & Truth optimal given beliefs & Weakest \\
\hline
\end{tabular}
\end{center}

%==============================================================================
% PART II: THE REVELATION PRINCIPLE
%==============================================================================
\subsection{Part II: The Revelation Principle – Truth as Benchmark}

\subsubsection{Direct Mechanisms}

A \textbf{direct mechanism} asks agents to report types:
\begin{equation}
S_i = \Theta_i \quad \text{(report your type)}
\end{equation}

\subsubsection{Incentive Compatibility}

A direct mechanism is \textbf{incentive compatible} if truth-telling is optimal:
\begin{equation}
\boxed{
U_i(\theta_i, \theta_i) \geq U_i(\hat{\theta}_i, \theta_i) \quad \forall \theta_i, \hat{\theta}_i
}
\label{eq:ic}
\end{equation}

Reporting true type $\theta_i$ is (weakly) better than any lie $\hat{\theta}_i$.

\subsubsection{The Revelation Principle}

Myerson (1979) proved the fundamental result:

\begin{equation}
\boxed{
\text{Any implementable } f \text{ can be implemented by incentive compatible direct mechanism}
}
\label{eq:revelation}
\end{equation}

\paragraph{Implication.}
To find all implementable outcomes, we can restrict attention to direct 
mechanisms where truth-telling is equilibrium. Simplifies analysis enormously.

\subsubsection{Framework Translation}

The Revelation Principle says:
\begin{equation}
C^*_{any mechanism} = C^*_{direct, truthful} \quad \text{(in terms of outcomes)}
\end{equation}

For design purposes, we can focus on $\Psi$ where truth is $C^*$.

%==============================================================================
% PART III: IMPOSSIBILITY RESULTS
%==============================================================================
\subsection{Part III: Impossibility Results – Limits of Design}

\subsubsection{Gibbard-Satterthwaite Theorem}

Gibbard (1973) and Satterthwaite (1975) proved:

\begin{equation}
\boxed{
\text{No strategy-proof, non-dictatorial mechanism exists for } \geq 3 \text{ alternatives}
}
\label{eq:gibbard-satterthwaite}
\end{equation}

\paragraph{Conditions.}
\begin{itemize}
    \item At least 3 alternatives
    \item All preference orderings possible (universal domain)
    \item Mechanism is onto (all alternatives can win)
\end{itemize}

Then: Either dictatorial or manipulable.

\subsubsection{Framework Translation}

Fundamental limit on $\Psi_{voting}$ design:
\begin{equation}
\nexists \Psi_{voting}: C^*_{truthful} \text{ and non-dictatorial and onto}
\end{equation}

\paragraph{Escape Routes.}
\begin{itemize}
    \item Restrict domain (single-peaked preferences)
    \item Allow randomization
    \item Relax strategy-proofness to Bayesian
    \item Accept some manipulation
\end{itemize}

\subsubsection{Myerson-Satterthwaite Impossibility}

Myerson \& Satterthwaite (1983) proved:

\begin{equation}
\boxed{
\text{No mechanism achieves efficiency + IR + BB for bilateral trade}
}
\label{eq:myerson-satterthwaite}
\end{equation}

where IR = individual rationality, BB = budget balance.

\paragraph{The Setting.}
\begin{itemize}
    \item Buyer values good at $v \sim F_v$
    \item Seller values at $c \sim F_c$
    \item Efficient: Trade iff $v > c$
    \item IR: Both willing to participate
    \item BB: No outside subsidy
\end{itemize}

\textbf{Result}: Cannot achieve all three simultaneously.

\subsubsection{Framework Translation}

Private information ($\Psi_K^{private}$) creates fundamental inefficiency:
\begin{equation}
\Psi_K^{asymmetric} \Rightarrow C^*_{feasible} \neq C^*_{first-best}
\end{equation}

Information rents are unavoidable cost of private information.

%==============================================================================
% PART IV: OPTIMAL AUCTION DESIGN
%==============================================================================
\subsection{Part IV: Optimal Mechanisms – Myerson's Auction}

\subsubsection{The Auction Design Problem}

Seller has one object, $n$ buyers with values $v_i \sim F_i$.
\begin{equation}
\max_{\text{mechanism}} \mathbb{E}[\text{Revenue}]
\end{equation}

\subsubsection{Myerson's Optimal Auction}

Myerson (1981) solved this completely:

\begin{equation}
\boxed{
\text{Allocate to highest } \psi_i(v_i) = v_i - \frac{1 - F_i(v_i)}{f_i(v_i)}
}
\label{eq:myerson-auction}
\end{equation}

where $\psi_i(v_i)$ is the \textbf{virtual valuation}.

\paragraph{Interpretation.}
\begin{itemize}
    \item Virtual valuation = value minus information rent
    \item Sell to highest virtual value (not highest value!)
    \item May exclude low-value buyers (reserve price)
\end{itemize}

\subsubsection{Revenue Equivalence}

Under certain conditions, all mechanisms with same allocation yield same revenue:

\begin{equation}
\boxed{
\text{Rev}(M_1) = \text{Rev}(M_2) \text{ if allocations identical}
}
\label{eq:revenue-equivalence}
\end{equation}

\paragraph{Implication.}
First-price, second-price, English, Dutch auctions all yield same expected 
revenue (with risk-neutral, IPV bidders).

\subsubsection{Framework Translation}

Optimal $\Psi_{auction}$ trades off efficiency and revenue:
\begin{equation}
\Psi^*_{auction} = \arg\max \text{Revenue} \quad \text{s.t. IC, IR}
\end{equation}

Information rents ($\Psi_K$) constrain what's achievable.

%==============================================================================
% PART V: VCG MECHANISMS
%==============================================================================
\subsection{Part V: VCG Mechanisms – Dominant Strategy Efficiency}

\subsubsection{Vickrey Auction}

Vickrey (1961) introduced the second-price sealed-bid auction:

\begin{equation}
\boxed{
\text{Winner pays second-highest bid}
}
\label{eq:vickrey}
\end{equation}

\paragraph{Key Property.}
Truth-telling is \textbf{dominant strategy}:
\begin{equation}
b_i^* = v_i \quad \forall v_i, \forall b_{-i}
\end{equation}

Optimal to bid true value regardless of what others do.

\subsubsection{Clarke-Groves Mechanism}

Clarke (1971) and Groves (1973) generalized to public goods:

\begin{equation}
\boxed{
t_i = \sum_{j \neq i} v_j(x^*) - \sum_{j \neq i} v_j(x^*_{-i})
}
\label{eq:vcg}
\end{equation}

Agent $i$ pays the externality imposed on others.

\paragraph{VCG Properties.}
\begin{itemize}
    \item Efficient: Chooses welfare-maximizing outcome
    \item Incentive compatible: Truth is dominant strategy
    \item Individually rational: Participation is optimal
    \item \textbf{Not} budget balanced (generally runs surplus)
\end{itemize}

\subsubsection{Framework Translation}

VCG is $\Psi_{mechanism}$ achieving efficiency via dominant strategies:
\begin{equation}
\Psi_{VCG} \Rightarrow C^*_{truthful} \Rightarrow X^*_{efficient}
\end{equation}

\paragraph{Limitation.}
Budget imbalance: $\sum_i t_i \neq 0$ generally.

\subsubsection{Green-Laffont Impossibility}

Green \& Laffont (1979) showed:

\begin{equation}
\boxed{
\text{No mechanism is efficient + dominant strategy IC + budget balanced}
}
\label{eq:green-laffont}
\end{equation}

Must sacrifice one of the three.

%==============================================================================
% PART VI: CONTRACT THEORY
%==============================================================================
\subsection{Part VI: Contract Theory – Mechanisms as Contracts}

\subsubsection{Moral Hazard}

Holmström (1979) analyzed optimal contracts under moral hazard:

\begin{equation}
\boxed{
\max_{w(x)} \mathbb{E}[x - w(x)] \quad \text{s.t. IC, IR}
}
\label{eq:moral-hazard}
\end{equation}

where $x$ is output, $w(x)$ is wage, agent chooses hidden effort $e$.

\paragraph{Key Results.}
\begin{itemize}
    \item Optimal contract depends on output (pay for performance)
    \item Trade-off: Incentives vs. risk-sharing
    \item First-best unattainable when effort unobservable
\end{itemize}

\subsubsection{Informativeness Principle}

Holmström proved:

\begin{equation}
\boxed{
\text{Optimal contract uses signal } y \text{ iff } y \text{ informative about effort}
}
\label{eq:informativeness}
\end{equation}

Any statistic correlated with effort should affect pay.

\subsubsection{Moral Hazard in Teams}

Holmström (1982) showed teams face additional problem:

\begin{equation}
\text{Budget balance} + \text{Nash IC} \Rightarrow \text{Inefficient effort}
\end{equation}

Cannot achieve efficient effort without budget-breaker.

\subsubsection{Framework Translation}

Contracts are $\Psi_{mechanism}$ for bilateral relationships:
\begin{equation}
\Psi_{contract} = w(x) \quad \text{(wage as function of output)}
\end{equation}

Optimal $\Psi_{contract}$ balances incentives and risk.

%==============================================================================
% PART VII: ROBUST MECHANISM DESIGN
%==============================================================================
\subsection{Part VII: Robust Mechanism Design – The Wilson Doctrine}

\subsubsection{The Problem with Common Knowledge}

Standard mechanism design assumes:
\begin{itemize}
    \item Common knowledge of type distributions
    \item Common knowledge of rationality
    \item Bayesian equilibrium
\end{itemize}

These are strong assumptions.

\subsubsection{Wilson Doctrine}

Robert Wilson advocated:

\begin{equation}
\boxed{
\text{Design mechanisms that don't rely on common knowledge assumptions}
}
\label{eq:wilson}
\end{equation}

\paragraph{Motivation.}
\begin{itemize}
    \item Where do common priors come from?
    \item What if agents have different models?
    \item Need mechanisms robust to misspecification
\end{itemize}

\subsubsection{Bergemann \& Morris: Robust Implementation}

Bergemann \& Morris (2005) formalized robust mechanism design:

\begin{equation}
\boxed{
\text{Implement } f \text{ for all possible beliefs, not just common prior}
}
\label{eq:robust}
\end{equation}

\paragraph{Key Results.}
\begin{itemize}
    \item Ex-post implementation (works for all type realizations)
    \item Stricter than Bayesian, but more robust
    \item VCG mechanisms are often robust
\end{itemize}

\subsubsection{Framework Translation}

Robust design accounts for $\Psi_K$ uncertainty:
\begin{equation}
\Psi^*_{robust} = \arg\max_{\Psi} \min_{\Psi_K} Q(C^*(\Psi, \Psi_K))
\end{equation}

Maximin over information structures.

%==============================================================================
% PART VIII: SYNTHESIS
%==============================================================================
\subsection{Part VIII: Synthesis – Mechanism Design Equations}

\subsubsection{Complete Framework Translation}

\begin{center}
\renewcommand{\arraystretch}{1.5}
\begin{tabular}{|l|c|l|}
\hline
\textbf{Concept} & \textbf{Equation} & \textbf{Framework Element} \\
\hline
Mechanism & $M = (S, g)$ & $\Psi_{mechanism}$ \\
\hline
Implementation & $g(C^*) = f(\theta)$ & $\Psi \to C^* \to X_{desired}$ \\
\hline
Revelation Principle & Restrict to direct IC & Truth as $C^*$ \\
\hline
Gibbard-Satterthwaite & No non-dictatorial SP & Limits on $\Psi_{voting}$ \\
\hline
Myerson-Satterthwaite & Eff + IR + BB impossible & $\Psi_K$ creates inefficiency \\
\hline
Myerson auction & Allocate by $\psi(v)$ & Optimal $\Psi_{auction}$ \\
\hline
VCG & Pay externality & Dominant strategy efficiency \\
\hline
Moral hazard & $w^*(x)$ & $\Psi_{contract}$ design \\
\hline
Robust design & Works for all $\Psi_K$ & Maximin over beliefs \\
\hline
\end{tabular}
\end{center}

\subsubsection{The Unified Picture}

Mechanism design shows:
\begin{equation}
\boxed{
C^*_{achievable} = f(\Psi_{mechanism}, \Psi_K, \text{constraints})
}
\end{equation}

What outcomes are achievable depends on:
\begin{enumerate}
    \item \textbf{Information structure} ($\Psi_K$): What's private?
    \item \textbf{Solution concept}: Dominant strategy, Nash, Bayesian?
    \item \textbf{Constraints}: IR, BB, etc.
\end{enumerate}

\subsubsection{Mechanism Design's Contribution to Framework}

\textbf{Without mechanism design:}
\begin{itemize}
    \item $\Psi$ taken as given
    \item No theory of institutional design
    \item No understanding of implementation limits
\end{itemize}

\textbf{With mechanism design:}
\begin{itemize}
    \item $\Psi$ as choice variable
    \item Theory of achievable $C^*$
    \item Impossibility results as design constraints
    \item Optimal mechanisms for specific goals
\end{itemize}

%==============================================================================
% PART IX: COMPLETE PAPER DOCUMENTATION
%==============================================================================
\subsection{Part IX: Complete Paper Documentation}

%--- FOUNDATIONS ---
\subsubsection{Foundations}

\paragraph{Paper 1: Hurwicz (1960)}

``Optimality and Informational Efficiency in Resource Allocation Processes.''
In \textit{Mathematical Methods in the Social Sciences}, Arrow, Karlin \& Suppes (eds.).

\textbf{Summary.}
Founding paper of mechanism design. Defines mechanisms, incentive compatibility, 
and asks what social choice functions can be implemented.

\textbf{Framework Translation.}
\begin{equation}
\Psi_{mechanism} = (S, g) \to C^* \to X
\end{equation}

\textbf{Key Finding.}
Nobel Prize (2007). Founded the field.

\paragraph{Paper 2: Hurwicz (1972)}

``On Informationally Decentralized Systems.''
In \textit{Decision and Organization}, Radner \& McGuire (eds.).

\textbf{Summary.}
Formalizes decentralized mechanisms. Shows competitive equilibrium can be 
implemented by decentralized mechanism.

\textbf{Framework Translation.}
Markets as mechanism: $\Psi_{market} \to C^*_{competitive}$.

%--- IMPOSSIBILITY ---
\subsubsection{Impossibility Results}

\paragraph{Paper 3: Gibbard (1973)}

``Manipulation of Voting Schemes: A General Result.''
\textit{Econometrica}, 41(4), 587--601.

\textbf{Summary.}
Proves no non-dictatorial, strategy-proof voting scheme exists with $\geq 3$ 
alternatives. Manipulation is always possible.

\textbf{Framework Translation.}
\begin{equation}
\nexists \Psi_{voting}: \text{Non-dictatorial and strategy-proof}
\end{equation}

\textbf{Key Finding.}
Fundamental impossibility for voting mechanisms.

\paragraph{Paper 4: Satterthwaite (1975)}

``Strategy-Proofness and Arrow's Conditions: Existence and Correspondence 
Theorems for Voting Procedures and Social Welfare Functions.''
\textit{Journal of Economic Theory}, 10(2), 187--217.

\textbf{Summary.}
Independent proof of Gibbard's theorem. Shows connection to Arrow's 
impossibility theorem.

\textbf{Framework Translation.}
Gibbard-Satterthwaite is mechanism design version of Arrow.

\paragraph{Paper 5: Myerson \& Satterthwaite (1983)}

``Efficient Mechanisms for Bilateral Trading.''
\textit{Journal of Economic Theory}, 29(2), 265--281.

\textbf{Summary.}
Proves impossibility of efficient bilateral trade with private values.
Cannot have efficiency + IR + budget balance.

\textbf{Framework Translation.}
\begin{equation}
\Psi_K^{private} \Rightarrow C^* \neq C^*_{efficient}
\end{equation}

\textbf{Key Finding.}
Fundamental limit on trade mechanisms. 3,000+ citations.

%--- VCG ---
\subsubsection{VCG Mechanisms}

\paragraph{Paper 6: Vickrey (1961)}

``Counterspeculation, Auctions, and Competitive Sealed Tenders.''
\textit{Journal of Finance}, 16(1), 8--37.

\textbf{Summary.}
Introduces second-price auction. Truth-telling is dominant strategy.
Also proposes VCG for multiple goods.

\textbf{Framework Translation.}
\begin{equation}
\Psi_{Vickrey}: C^*_{truthful} \text{ in dominant strategies}
\end{equation}

\textbf{Key Finding.}
Nobel Prize (1996). Foundation of auction theory.

\paragraph{Paper 7: Clarke (1971)}

``Multipart Pricing of Public Goods.''
\textit{Public Choice}, 11(1), 17--33.

\textbf{Summary.}
Extends Vickrey to public goods. Agents pay their externality on others.

\textbf{Framework Translation.}
Clarke mechanism for public goods provision.

\paragraph{Paper 8: Groves (1973)}

``Incentives in Teams.''
\textit{Econometrica}, 41(4), 617--631.

\textbf{Summary.}
General class of mechanisms with dominant strategy truth-telling.
Groves mechanisms include Vickrey and Clarke as special cases.

\textbf{Framework Translation.}
\begin{equation}
t_i = h_i(\theta_{-i}) + \sum_{j \neq i} v_j(x(\theta), \theta_j)
\end{equation}

\paragraph{Paper 9: Groves \& Ledyard (1977)}

``Optimal Allocation of Public Goods: A Solution to the `Free Rider' Problem.''
\textit{Econometrica}, 45(4), 783--809.

\textbf{Summary.}
Mechanism for efficient public goods with Nash implementation.
Not dominant strategy, but achieves efficiency and budget balance.

\textbf{Framework Translation.}
Trade-off: Dominant strategy vs. budget balance.

%--- OPTIMAL MECHANISMS ---
\subsubsection{Optimal Mechanism Design}

\paragraph{Paper 10: Myerson (1979)}

``Incentive Compatibility and the Bargaining Problem.''
\textit{Econometrica}, 47(1), 61--73.

\textbf{Summary.}
Proves Revelation Principle. Any implementable outcome can be achieved 
by incentive compatible direct mechanism.

\textbf{Framework Translation.}
\begin{equation}
\text{Can restrict to } C^*_{truthful}
\end{equation}

\textbf{Key Finding.}
Foundational methodological result. 4,000+ citations.

\paragraph{Paper 11: Myerson (1981)}

``Optimal Auction Design.''
\textit{Mathematics of Operations Research}, 6(1), 58--73.

\textbf{Summary.}
Solves revenue-maximizing auction. Introduces virtual valuations.
Optimal to allocate by virtual value, not actual value.

\textbf{Framework Translation.}
\begin{equation}
\Psi^*_{auction} = \text{Allocate to highest } \psi_i(v_i)
\end{equation}

\textbf{Key Finding.}
Nobel Prize (2007). 6,000+ citations.

%--- IMPLEMENTATION THEORY ---
\subsubsection{Implementation Theory}

\paragraph{Paper 12: Maskin (1999)}

``Nash Equilibrium and Welfare Optimality.''
\textit{Review of Economic Studies}, 66(1), 23--38.

\textbf{Summary.}
Characterizes Nash-implementable social choice functions.
Monotonicity is necessary; with no veto power, sufficient.

\textbf{Framework Translation.}
\begin{equation}
f \text{ Nash-implementable} \Leftrightarrow f \text{ monotonic (+ NVP)}
\end{equation}

\textbf{Key Finding.}
Nobel Prize (2007). Circulated since 1977.

\paragraph{Paper 13: Maskin \& Sjöström (2002)}

``Implementation Theory.''
\textit{Handbook of Social Choice and Welfare}, Vol. 1, 237--288.

\textbf{Summary.}
Comprehensive survey of implementation theory. Covers Nash, dominant 
strategy, Bayesian, virtual implementation.

\textbf{Framework Translation.}
Complete theory of which $C^*$ are achievable under which solution concepts.

%--- CONTRACT THEORY ---
\subsubsection{Contract Theory}

\paragraph{Paper 14: Holmström (1979)}

``Moral Hazard and Observability.''
\textit{Bell Journal of Economics}, 10(1), 74--91.

\textbf{Summary.}
Optimal contracts under moral hazard. Proves informativeness principle:
use any signal correlated with effort.

\textbf{Framework Translation.}
\begin{equation}
\Psi^*_{contract} \text{ uses all informative signals}
\end{equation}

\textbf{Key Finding.}
Nobel Prize (2016). Foundation of contract theory.

\paragraph{Paper 15: Holmström (1982)}

``Moral Hazard in Teams.''
\textit{Bell Journal of Economics}, 13(2), 324--340.

\textbf{Summary.}
Teams face budget balance constraint. Cannot achieve efficient effort 
with Nash IC and budget balance.

\textbf{Framework Translation.}
Teams need budget-breaker for efficiency.

\textbf{Key Finding.}
2,500+ citations.

%--- AUCTIONS ---
\subsubsection{Auction Theory}

\paragraph{Paper 16: Milgrom \& Weber (1982)}

``A Theory of Auctions and Competitive Bidding.''
\textit{Econometrica}, 50(5), 1089--1122.

\textbf{Summary.}
General auction theory with affiliated values. English auction outperforms 
sealed-bid when values affiliated.

\textbf{Framework Translation.}
Optimal $\Psi_{auction}$ depends on information structure.

\textbf{Key Finding.}
Nobel Prize (2020). 5,000+ citations.

\paragraph{Paper 17: Wilson (1987)}

``Game-Theoretic Analysis of Trading Processes.''
In \textit{Advances in Economic Theory}, Bewley (ed.).

\textbf{Summary.}
Advocates ``Wilson Doctrine'': design mechanisms that don't rely on 
common knowledge of distributions.

\textbf{Framework Translation.}
Robust $\Psi$ design.

\textbf{Key Finding.}
Nobel Prize (2020).

%--- ROBUST DESIGN ---
\subsubsection{Robust Mechanism Design}

\paragraph{Paper 18: Bergemann \& Morris (2005)}

``Robust Mechanism Design.''
\textit{Econometrica}, 73(6), 1771--1813.

\textbf{Summary.}
Formalizes robust implementation. Mechanism works for all possible 
beliefs, not just common prior.

\textbf{Framework Translation.}
\begin{equation}
\Psi^*_{robust} \text{ works } \forall \Psi_K
\end{equation}

\textbf{Key Finding.}
Major advance in mechanism design theory.

%--- TEXTBOOKS ---
\subsubsection{Textbooks and Surveys}

\paragraph{Paper 19: Klemperer (2004)}

\textit{Auctions: Theory and Practice.}
Princeton University Press.

\textbf{Summary.}
Accessible treatment of auction theory and design. Covers theory, 
practical design, and case studies (3G spectrum auctions).

\textbf{Framework Translation.}
Applied $\Psi_{auction}$ design.

\paragraph{Paper 20: Börgers (2015)}

\textit{An Introduction to the Theory of Mechanism Design.}
Oxford University Press.

\textbf{Summary.}
Modern graduate textbook. Covers foundations, implementation, auctions, 
bilateral trade, public goods, and robust design.

\textbf{Framework Translation.}
Complete theory of $\Psi$ design.

%%%%%%%%%%%%%%%%%%%%%%%%%%%%%%%%%%%%%%%%%%%%%%%%%%%%%%%%%%%%%%%%%%%%%%%%%%%%%%%
\subsection{Citation and Verification Note}

\textbf{Nobel Prizes represented:}
\begin{itemize}
    \item Vickrey (1996) ``for contributions to the economic theory of incentives under asymmetric information''
    \item Hurwicz, Maskin \& Myerson (2007) ``for having laid the foundations of mechanism design theory''
    \item Holmström (2016) ``for contributions to contract theory''
    \item Milgrom \& Wilson (2020) ``for improvements to auction theory and inventions of new auction formats''
\end{itemize}

Combined Google Scholar citations for the 20 papers: $>$60,000.

This appendix provides the theoretical foundations for institutional design 
within the $(C, \Psi, C^*, K, Q)$ framework, showing what outcomes are 
achievable through mechanism design and what fundamental limits exist.

%%%%%%%%%%%%%%%%%%%%%%%%%%%%%%%%%%%%%%%%%%%%%%%%%%%%%%%%%%%%%%%%%%%%%%%%%%%%%%%
